\documentclass[10.5pt,a4paper]{article}

% ---------- Packages ----------
\usepackage[utf8]{inputenc}
\usepackage[T1]{fontenc}
\usepackage[french]{babel}
\usepackage{XCharter}                % elegant serif body font
\usepackage[a4paper,margin=1.1cm]{geometry}
\usepackage{xcolor}
\usepackage{titlesec}
\usepackage{enumitem}
\usepackage{hyperref}
\usepackage{ragged2e}
\usepackage{paracol}
\usepackage[none]{hyphenat}
\setlength{\emergencystretch}{3em}

% ---------- Colors ----------
\definecolor{accent}{RGB}{30,64,149}
\definecolor{graytext}{RGB}{90,90,90}

% ---------- Links ----------
\hypersetup{
  colorlinks=true,
  linkcolor=accent,
  urlcolor=accent
}

% ---------- Section titles (main column) ----------
\titleformat{\section}
  {\color{accent}\bfseries\Large}
  {}{0em}{}
\titlespacing{\section}{0pt}{6pt}{3pt}
\newcommand{\sectionrule}{\vspace{-6pt}{\color{accent}\hrule height 1pt}\vspace{3pt}}

% ---------- Sidebar section titles ----------
\newcommand{\subsectitle}[1]{%
  {\color{accent}\bfseries\large #1}\\[-2pt]
  {\color{accent}\hrule height 0.8pt}\vspace{5pt}
}

\setlength{\parindent}{0pt}
\pagestyle{empty}

% ---------- Custom commands ----------

% Experience / generic entry: #1 position #2 org #3 location #4 period #5 body
\newcommand{\cvexp}[5]{%
  \noindent\textbf{#1} \hfill {\itshape\color{graytext}\small #4}\\
  {\itshape #2 --- #3}\\[-2pt]
  #5
  \vspace{4pt}
}

% Education entry: #1 school #2 area/degree #3 location #4 period #5 description
\newcommand{\cvedu}[5]{%
  \noindent\textbf{#1} \hfill {\itshape\color{graytext}\small #4}\\
  {\itshape #2, #3}\\[1pt]
  {\small #5}
  \vspace{3pt}
}

% Education entry without description
\newcommand{\cveduplain}[4]{%
  \noindent\textbf{#1} \hfill {\itshape\color{graytext}\small #4}\\
  {\itshape #2, #3}
  \vspace{3pt}
}

% Project entry: #1 name #2 period #3 description
\newcommand{\cvproject}[3]{%
  \noindent\textbf{#1} \hfill {\itshape\color{graytext}\small #2}\\[1pt]
  {\small #3}
  \vspace{3pt}
}

% Sidebar simple line entry: #1 title (bold) #2 detail
\newcommand{\sidebarline}[2]{%
  \noindent\textbf{\small #1}\\
  {\small #2}\\[3pt]
}

% Sidebar dated entry (volunteer, with description)
\newcommand{\sidebardated}[3]{%
  \textbf{\small #1} {\itshape\color{graytext}\footnotesize (#2)}\\
  {\small #3}\\[3pt]
}

% Sidebar dated entry without description (certifications)
\newcommand{\sidebarcert}[2]{%
  \noindent\textbf{\small #1} {\itshape\color{graytext}\footnotesize (#2)}\\[2pt]
}

\begin{document}

% ================= HEADER =================
\begin{center}
  {\fontsize{26}{30}\selectfont\bfseries Amélie Yang}\\[4pt]
  {\color{accent}\itshape\large Étudiante ingénieure en Data \& Systèmes d'Information -- IMT Mines Albi}\\[6pt]
  {\small
    Grenoble, France
    \,\textbullet\,
    \href{mailto:amelie.yang@mines-albi.fr}{amelie.yang@mines-albi.fr}
    \,\textbullet\,
    +33~7~67~95~35~99
    \,\textbullet\,
    \href{https://www.linkedin.com/in/amelieyang}{linkedin.com/in/amelieyang}
  }
\end{center}

\vspace{3pt}
{\color{accent}\hrule height 1.2pt}
\vspace{4pt}

% ================= TWO COLUMN BODY =================
\columnratio{0.65}
\setlength{\columnsep}{1.5em}
\begin{paracol}{2}

% ---------- RÉSUMÉ ----------
\section*{Résumé}
\sectionrule
\justifying
Ingénieure diplômée de l'IMT Mines Albi (Data \& Systèmes d'Information), avec une solide formation en mathématiques appliquées et une expérience en développement d'algorithmes de machine learning, optimisation et traitement de données en Python. Rigueur expérimentale (protocoles de test reproductibles, métriques quantifiées) et intérêt pour l'application de l'IA à de nouveaux domaines techniques. En recherche d'un poste d'ingénieure-chercheuse à partir de septembre 2026.

% ---------- EXPÉRIENCE ----------
\section*{Expérience}
\sectionrule

\cvexp{Stagiaire Ingénieure IA Générative}{ALTEN Lab}{Grenoble, France}{Février 2026 -- Août 2026 (6 mois)}{%
  \begin{itemize}[leftmargin=12pt, itemsep=2pt, topsep=0pt, partopsep=0pt, parsep=0pt]
    \item Conception et validation expérimentale d'algorithmes de traitement de données en Python, avec protocole de test reproductible et métriques quantifiées (WER, CER, NED, GriTS, F1).
    \item Développement de deux architectures algorithmiques complémentaires (frameworkless, LangChain/\allowbreak LangGraph), comparées pour identifier la plus performante.
    \item Intégration et validation du système sur une plateforme conteneurisée (API FastAPI, Docker), avec déploiement en production (CI/CD).
  \end{itemize}
}

\cvexp{Stagiaire Assistante Ingénieure R\&D}{Yooz}{Aimargues, France}{Avril 2025 -- Août 2025 (4 mois)}{%
  \begin{itemize}[leftmargin=12pt, itemsep=2pt, topsep=0pt, partopsep=0pt, parsep=0pt]
    \item Développement et comparaison de modèles de machine learning (Random Forest, XGBoost) sur données réelles bruitées.
    \item Analyse de l'espace vectoriel par réduction de dimension (PCA, t-SNE, UMAP) pour diagnostiquer les limites d'un système en production.
    \item Étude comparative quantifiée de plusieurs approches algorithmiques (81\% vs 64\% de performance), pour orienter le choix de solution selon les contraintes du contexte.
  \end{itemize}
}

\cvexp{Stagiaire Opératrice}{Laboratoires Pierre Fabre}{Soual, France}{Février 2024 -- Mars 2024 (2 mois)}{%
  \begin{itemize}[leftmargin=12pt, itemsep=2pt, topsep=0pt, partopsep=0pt, parsep=0pt]
    \item Renforcement des compétences opérationnelles en gérant la ligne de conditionnement, surveillance, tri et alimentation des processus.
  \end{itemize}
}

% ---------- FORMATION ----------
\section*{Formation}
\sectionrule

\cvedu{The Catholic University of Korea}{AI Major -- Exchange Student}{Séoul, Corée du Sud}{Septembre 2025 -- Décembre 2025}{Computer Vision, Data Mining, Machine Learning.}

\cvedu{IMT Mines Albi}{Ingénieur Généraliste}{Albi, France}{2023 -- 2026}{Optimisation, Base de données, Calcul numérique, Big Data, Machine Learning.}

\cvedu{Lycée Joffre}{Classe préparatoire MPSI-MP}{Montpellier, France}{2021 -- 2023}{Mathématiques appliquées, physique, algorithmique.}

% ---------- PROJETS ----------
\section*{Projets}
\sectionrule

\cvproject{Prévision en temps réel du temps d'attente à quai de la SNCF}{2025}{Nettoyage, traitement et agrégation de données de flux bruitées et variables. Entraînement et évaluation de modèles de régression et de séries temporelles en Python pour la prédiction en temps réel.}

\switchcolumn

\vspace{2pt}

% ---------- COMPÉTENCES ----------
\subsectitle{Compétences}

\sidebarline{Mathématiques \& Algorithmique}{Mathématiques appliquées, optimisation, calcul numérique}
\sidebarline{Langages}{Python (NumPy, Pandas, Scikit-learn, TensorFlow/\allowbreak PyTorch), Matlab, C++, R, SQL}
\sidebarline{Machine Learning \& IA}{CNN, MLP, régression, séries temporelles, classification, clustering}
\sidebarline{Outils}{Git, Docker, FastAPI, CI/CD}

% ---------- LANGUES ----------
\subsectitle{Langues}
\noindent{\small Anglais (TOEIC 965 pts), Espagnol, Japonais}

\vspace{5pt}

% ---------- ASSOCIATIF ----------
\subsectitle{Associatif}

\sidebardated{Secrétaire Générale -- Bureau des Sports IMT Mines Albi}{2024 -- 2025}{Organisation de plannings et d'événements. Rédaction de dossiers, comptes-rendus.}

\sidebardated{Les Cordées de la Réussite}{2023 -- 2024}{Encadrement et accompagnement d'élèves dans leur orientation.}

% ---------- CERTIFICATIONS ----------
\subsectitle{Certifications}

\sidebarcert{MOOC -- Prévention des risques professionnels (INRS)}{2023}
\sidebarcert{MOOC -- Propriété Intellectuelle (INPI)}{2024}

\vspace{4pt}

% ---------- INTÉRÊTS ----------
\subsectitle{Intérêts}
\noindent{\small Volley (2 ans en compétition)}\\
{\small Musique (piano, guitare)}\\
{\small Jeux-vidéos}

\end{paracol}

\end{document}
